\begin{textsample}{POS Dim 4 – vem\_ed\_35 – Score 97.00 – vem\_ed\_35/t187.txt}  \label{ex:f4_pos_001}
5º IVES 2026: obstáculos e pontes Com o tema “Cooperation in COIL Virtual Exchange: hurdles and bridges” (Cooperação em Intercâmbios Virtuais do \textbf{tipo} COIL: obstáculos e pontes), a quinta edição do International Virtual Exchange Symposium (IVES CGESG 2026) reuniu quatro palestrantes internacionais: Sara Pittarello (UNICollaboration, rede \textbf{europeia} de apoio aos Intercâmbios Virtuais), Carlos Maza (UNAM, México), Henry Shepherd (IVEC, EUA) e Salome Escalona (Bukidnon State University, Filipinas).
O evento online ocorreu em 28 de \textbf{maio} de 2026 e engajou 189 participantes.
Sara Pittarello, diretora da UNICollaboration, ressaltou que os Intercâmbios Virtuais permitem expandir acesso a experiências internacionais, desenvolver competências globais e habilidades interculturais, \textbf{apoiam} equidade e inclusão e preparam os estudantes para ambientes de trabalho \textbf{globalizados}.
Não se trata apenas de \textbf{conectividade}, mas de interações significativas, objetivos compartilhados, colaboração estruturada, aprendizagem \textbf{cocriada} e parcerias de longo \textbf{prazo}.
Apresentou projetos financiados pelo Erasmus+, principal programa de mobilidade acadêmica \textbf{europeia}, que abordam temas como “Critical Virtual Exchange in AI”, voltado para uma pedagogia crítica, a consciência ética e o letramento em Inteligência Artificial por meio de Intercâmbios Virtuais entre Europa e África subsaariana; desenvolvimento profissional de professores com tecnologias digitais e o programa VAMOS, voltado a Intercâmbios Virtuais que \textbf{solucionam} problemas de desenvolvimento sustentável.
O VAMOS integrou universidades da Europa, do Brasil e de Honduras.
Para encerrar, Pittarello destacou quatro pontos: - Intercâmbios Virtuais são mais que tecnologia - Sustentabilidade dos projetos requer estrutura e estratégia - Facilitação é fundamental, não opcional - As parcerias devem ser de propriedade conjunta Carlos Maza, da UNAM, destacou a construção de “pontes com \textbf{alicerces} humanos” por meio da abordagem COIL (Collaborative Online International Learning).
Lembrou como a pandemia de covid-19 impulsionou rapidamente a comunidade \textbf{universitária} para o ensino \textbf{remoto}, criando as bases para a colaboração internacional virtual.
Com a retomada das atividades \textbf{presenciais}, a UNAM \textbf{aproveitou} essa infraestrutura para implementar o programa COIL, conectando salas de aula locais a instituições ao redor do mundo.
Nesse contexto, a \textbf{adesão} a redes como COIL Connect e Red LatAM COIL foi essencial para estruturar parcerias globais.
A experiência demonstrou que o principal pilar são os docentes.
Para sua formação continuada, a universidade oferece \textbf{oficinas} de capacitação e eventos acadêmicos como as Jornadas COIL UNAM.
Além de formar professores, é importante preparar os alunos para as experiências internacionais virtuais, reduzindo \textbf{ansiedade} e promovendo engajamento \textbf{consciente}.
Maza destaca que o COIL democratiza o acesso à internacionalização, ampliando oportunidades em uma universidade grande como a UNAM (que tem mais de 370 mil alunos).
Entre os desafios, destaca a \textbf{difusão} do conceito COIL para toda a comunidade acadêmica e a consolidação de diretrizes institucionais.
Henry Shepherd, do International Virtual Exchange Consortium (IVEC), destacou dois obstáculos aos projetos COIL e duas pontes para superá-los.
O primeiro é a dificuldade de reconhecer diferenças de perspectiva entre parceiros. “Até mesmo a \textbf{palavra} COIL não \textbf{significa} a mesma coisa em todos os lugares”, observou.
Nossas experiências e preferências moldam nossas suposições sobre cada aspecto do COIL”. \textbf{Deixar} de considerar o ponto de \textbf{vista} do outro é \textbf{fonte} de muitos desafios nos projetos.
Para superar esse obstáculo, Shepherd sugere colocar-se no lugar do parceiro, tomar \textbf{decisões} conjuntamente desde o planejamento \textbf{inicial} e, quando não houver meio-termo entre os parceiros, revezar: “sua plataforma, meu app; horário da sua aula primeiro, depois o meu”.
O segundo obstáculo é a \textbf{incerteza}.
Geralmente, apontamos os desafios óbvios: diferenças de fusos horários, de tecnologias, \textbf{calendários} ocupados.
O verdadeiro problema pode estar em outro ponto: diferenças nos \textbf{estilos} de comunicação, de gestão de projetos, ou de abordar uma questão que não foi tratada. “A comunicação que para uns \textbf{parece} profissional, para outros pode ser sentida como fria.
É difícil \textbf{ler} comportamentos \textbf{online} em alguém com quem você nunca encontrou pessoalmente”, \textbf{ponderou}.
Além disso, as culturas expressam o \textbf{atrito} de formas diferentes. “Podemos simplesmente tentar superar isso e seguir em frente, podemos nos afastar e nos distanciar do nosso parceiro, ou ir embora sem dar explicações e \textbf{deixar} o projeto ir por \textbf{água} abaixo.
Infelizmente, tudo isso acontece no COIL”.
Por isso, Shepherd \textbf{encoraja} os parceiros a conversarem sobre expectativas sobre como vão trabalhar \textbf{juntos} e \textbf{encarar} os desafios, antes mesmo que aconteçam.
Reconhecer que as pessoas \textbf{colaboram} de formas diferentes: silêncio não \textbf{significa} desinteresse.
E \textbf{resolver} \textbf{atritos} antes que se prolonguem por semanas.
Outra recomendação é incluir nas \textbf{reuniões} de planejamento com o professor parceiro e na preparação para o projeto com os alunos, uma discussão sobre a complexidade e as oportunidades de aprendizagem que \textbf{surgem} ao \textbf{lidar} com a cooperação e a colaboração a \textbf{distância}.
Para que os docentes possam aprender com os \textbf{pares}, fortalecendo práticas, compartilhando experiências e soluções, Shepherd \textbf{recomenda} o engajamento em redes internacionais como COIL Connect, Red LatAM COIL, BraVE, UNICollaboration, AfriCOIL, Japan COIL Association. “Quando \textbf{usamos} a paciência e a \textbf{curiosidade} para compreender melhor os pontos de \textbf{vista} dos outros, construímos relacionamentos, \textbf{desenvolvemos} habilidades críticas úteis e encontramos satisfação por ter feito algo difícil, e isso realmente \textbf{vale} a \textbf{pena} fazer”.
Salome Escalona, da Bukidnon State University (BukSU, Filipinas), apresentou obstáculos e pontes sob o ponto de \textbf{vista} de coordenadores e professores.
Citou que, ao \textbf{ingressar} em ambientes globais sofisticados, repletos de profissionais altamente \textbf{experientes}, professores podem \textbf{experimentar} a “síndrome do impostor”.
Para superar essa síndrome, é importante valorizar a \textbf{mentoria} e a construção de relacionamentos, aprendendo com a generosidade de parceiros \textbf{experientes}.
Alguns professores podem se assustar com a \textbf{carga} de trabalho gerada pelos projetos COIL, que podem ser vistos como “tarefa \textbf{extra}”.
Para evitar essa visão, coordenadores de projetos COIL devem demonstrar como as atividades podem ser integradas aos objetivos de aprendizagem da disciplina.
Assim, os professores integram o projeto COIL a tarefas que já fazem parte do plano de ensino.
Como \textbf{lidar} com a \textbf{ansiedade} dos estudantes e barreiras de comunicação?
A professora sugere a construção de um ambiente de comunicação seguro e encorajador.
Para contornar a \textbf{redução} do engajamento \textbf{estudantil} e a fadiga digital, Salome Escalona sugere um ritmo de atividades mais espaçado e \textbf{flexível}, com tarefas ancoradas em conexões humanas reais. “Quando confrontados com a tarefa de produzir um trabalho verdadeiramente intercultural, os alunos \textbf{deixaram} de \textbf{encarar} a atividade como apenas mais uma tarefa”.
E conclui com uma reflexão: “O intercâmbio virtual é muito mais do que uma \textbf{exigência} \textbf{administrativa} ou uma \textbf{lista} de parceiros internacionais.
É um forte lembrete de que as diferenças \textbf{geográficas}, culturais e linguísticas \textbf{desaparecem} quando priorizamos a conexão humana”.
Os \textbf{quadros} a seguir \textbf{sintetizam} os obstáculos e pontes sob as perspectivas dos coordenadores e professores, apresentadas por Salome Escalona (\textbf{leia} sobre o PCI realizado entre Fatecs Pompeia, Guarulhos e Zona Leste e BukSU no “Artigo de Opinião”).
Histórico do IVES As quatro edições anteriores do International Virtual Exchange Symposium (IVES), realizadas entre 2022 e 2025, foram registradas nos números 15, 18, 24 e 29 de VEm e publicadas em \textbf{playlists} do Canal CGESG no YouTube: IVES 2025: https://bit.ly/3GWupxd IVES 2024: http://bit.ly/3SQtvFi IVES 2023: https://bit.ly/4k3ymPq IVES 2022: https://bit.ly/3FnE5R9 Ao longo dos últimos anos, a \textbf{trajetória} do IVES \textbf{vem} consolidando os Intercâmbios Virtuais como estratégia de Internacionalização em Casa para o Ensino Superior no Centro Paula Souza, bem como a evolução dos Projetos Colaborativos Internacionais (PCIs) realizados nas Fatecs sob a coordenação da equipe de Apoio à Internacionalização do Ensino Superior da Coordenadoria Geral de Ensino Superior de Graduação do CPS.
A primeira edição (IVES 2022) teve caráter \textbf{inaugural}, reunindo parceiros internacionais (China, Índia, África do Sul e Colômbia) e apresentando PCIs desenvolvidos nas Fatecs Jaboticabal, Garça, Franco da Rocha, São José do Rio Preto, Matão e Itu.
Os relatos de experiência apresentaram a abordagem COIL (Collaborative Online International Learning) tropicalizada nos PCIs.
Também foi mencionada a \textbf{criação} do ProCin nas Etecs, programa iniciado em 2021 pela Coordenadoria Geral de Ensino Médio e Técnico e pela ARInter, com a \textbf{mentoria} da equipe dos PCIs.
O evento destacou o Intercâmbio Virtual como ferramenta de Internacionalização em Casa, baseada na colaboração \textbf{remota} e no desenvolvimento de competências interculturais.
Na segunda edição (IVES 2023), houve a participação de cinco \textbf{especialistas} internacionais e ampliação do público, que \textbf{subiu} de 176 para 190 participantes.
As \textbf{palestras} foram \textbf{oferecidas} por coordenadores de Intercâmbios Virtuais da DePaul University (EUA), DUOC UC (Chile), Universidad Católica Andrés Bello (Venezuela), University of Minnesota System e SUNY COIL Center (EUA).
As discussões enfatizaram \textbf{tendências} e desafios do COIL.
Hope Windle (SUNY COIL Center) trouxe conceitos como inclusão, \textbf{interseccionalidade} e “esperança crítica”, além de \textbf{projeções} para o futuro dos intercâmbios virtuais.
O evento se posiciona como \textbf{espaço} de troca global de boas práticas.
A terceira edição (IVES 2024) teve público \textbf{recorde} (246 participantes) e \textbf{aprofundou} o debate ao propor o tema “Beyond Borders: Evolving COIL Virtual Exchange” (Além das Fronteiras: Intercâmbio Virtual COIL em evolução).
O foco \textbf{mudou} da apresentação de experiências para a análise de estratégias de implementação, \textbf{manutenção} e expansão dos projetos, com \textbf{destaque} para novas abordagens pedagógicas e formatos \textbf{inovadores} de interação.
Os palestrantes internacionais foram Mirjam Hauck (Open University, Reino Unido); Andrea Alsina e Sofia Pekarek (Universidad Católica de Salta, Argentina); Sandra Valencia, (Universidad Católica de Manizales, Colômbia); Brenda García (Red LatAM COIL, México) e Gabriela Méndez (Florida International University, EUA).
A quarta edição (IVES 2025) \textbf{reafirmou} o evento como fórum internacional de reflexão estratégica, com a temática Changing Horizons (Mudança de Horizontes).
As discussões enfatizaram o papel dos intercâmbios virtuais frente às transformações educacionais contemporâneas, reunindo \textbf{especialistas} de diferentes países: João Monteiro, do ISP Gaya (Portugal), Stephanie Doscher, da University of Minnesota (EUA), Julia Bacca e Diana Garzon, da Uniminuto (Colômbia) e Adam Freed, da University of Michigan (EUA).
O evento \textbf{evidenciou} a importância da troca cultural e da formação global dos estudantes, além de \textbf{dialogar} com agendas mais \textbf{amplas} de internacionalização.
Ao longo dos anos, o IVES evoluiu de um encontro voltado à apresentação institucional para um \textbf{espaço} consolidado de debate global sobre Intercâmbios Virtuais.
Observa-se uma maior complexidade das discussões, passando da \textbf{divulgação} de práticas ao debate sobre perspectivas estratégicas e inovação pedagógica.
Essa \textbf{trajetória} \textbf{reflete} a expansão dos PCIs e o fortalecimento do Centro Paula Souza como referência internacional em COIL, iniciativa essencial para a internacionalização \textbf{inclusiva} e colaborativa no ensino superior.

% matched lemmas: adesão, administrativo, alicerce, ampla, ansiedade, apoiar, aprofundar, aproveitar, atrito, calendário, carga, cocriar, colabor, conectividade, consciente, criação, curiosidade, decisão, deixar, desaparecer, desenvolvemos, destaque, dialogar, difusão, distância, divulgação, encarar, encoraja, espaço, especialista, estilo, estudantil, europeu, evidenciar, exigência, experiente, experimentar, extra, flexível, fonte, geográfico, globalizado, inaugural, incerteza, inclusivo, ingressar, inicial, inovador, interseccionalidade, junto, ler, lidar, lista, maio, manutenção, mentoria, mudar, oferecir, oficina, onlinar, palavra, palestra, par, parecer, pena, playlists, ponderar, prazo, presencial, projeção, quadro, reafirmar, recomendar, recorde, redução, refleter, remoto, resolver, reunião, significar, sintetizar, solucionar, subir, surgir, tendência, tipo, trajetória, universitário, usar, valer, vir, vista, água
\end{textsample}
